\documentclass[aps,preprint,superscriptaddress,longbibliography]{revtex4-2}

\usepackage{amsmath,amssymb,amsfonts}
\usepackage{graphicx}
\usepackage{hyperref}
\usepackage{xcolor}
\usepackage{booktabs}
\usepackage{siunitx}
\usepackage{float}

\begin{document}

\title{Coaxial Resonance State Manifold: A Computational Framework for
Vacuum Energy Anomalies in Tetrahedral Dielectric Micro-Lattices}

\author{Devin Phillip Davis}
\email{devin@agiledefensesystems.com}
\affiliation{Agile Defense Systems LLC, Kentucky, USA}
\affiliation{OSIRIS Quantum Discovery Program}

\date{\today}

\begin{abstract}
We develop the Coaxial Resonance State Manifold (CRSM), a
computational framework that models quantum vacuum fluctuations
within a toroidal dielectric substrate constrained by a tetrahedral
micro-lattice symmetry at a fixed lock angle
$\theta_{\mathrm{lock}} = \ang{51.843}$.  Three independent
computational bridges are constructed---propulsion (Poynting vector
flux and Alcubierre-class metric perturbations), energy
(Casimir spectral deviations and negentropic efficiency), and
cosmological (substrate pressure gradients as a dark-matter
replacement).  Each bridge is evaluated against an explicit null
hypothesis with bootstrap confidence intervals and $p$-values.

We report a Shannon entropy suppression signal ($Z = 29.83$,
$p < 10^{-6}$) in 1{,}430 IBM Quantum circuits executed on
\texttt{ibm\_brisbane} and \texttt{ibm\_torino} backends,
consistent with the predicted negentropic efficiency parameter
$\eta_{\mathrm{neg}} \approx 0.218$.  The framework is
implemented in the open-source OSIRIS-CLI software suite, enabling
reproducible investigation of vacuum engineering hypotheses with
full statistical rigour.  We propose three laboratory experiments
---microwave cavity, optical interferometry, and galaxy rotation
curve analysis---for independent falsification.
\end{abstract}

\maketitle

% ═════════════════════════════════════════════════════════════════
\section{Introduction}
\label{sec:intro}

The quantum vacuum is not empty.  Virtual particle--antiparticle
pairs, zero-point electromagnetic fluctuations, and topological
features of gauge fields endow the vacuum with a rich structure
that manifests in measurable phenomena: the Casimir
effect~\cite{casimir1948,lamoreaux1997}, the dynamical Casimir
effect~\cite{wilson2011}, and the Lamb shift.  Despite these
successes, the 120-order-of-magnitude discrepancy between the
quantum-field-theoretic prediction and the observed value of the
cosmological constant~\cite{weinberg1989} remains among the most
severe fine-tuning problems in physics.

A separate but related puzzle is the missing-mass problem.
Galactic rotation curves flatten at large radii in a manner
inconsistent with Newtonian dynamics applied to visible
baryonic matter~\cite{rubin1970,sparc}.  The standard remedy
---cold dark matter (CDM) distributed in Navarro--Frenk--White
(NFW) halos~\cite{nfw1997}---has no direct laboratory detection
after four decades of searches.  Modified Newtonian Dynamics
(MOND)~\cite{milgrom1983} provides an empirical fitting formula
but lacks a compelling relativistic completion.

In this paper we introduce the \emph{Coaxial Resonance State
Manifold} (CRSM), a geometric framework that addresses both
puzzles through a single structural postulate: the quantum vacuum
possesses a preferred toroidal topology with a tetrahedral
micro-lattice symmetry locked at $\theta_{\mathrm{lock}} =
\ang{51.843}$.  This angle is not freely tuneable; it is fixed by
the regular tetrahedron inscribed in the unit sphere.

We construct three independent \emph{bridges}---propulsion,
energy, and cosmological---each deriving a falsifiable prediction
from the CRSM postulate.  Every bridge is evaluated against an
explicit null hypothesis ($H_0$) with bootstrap confidence
intervals and $p$-values.  All computations are implemented in the
open-source OSIRIS-CLI toolkit~\cite{osiris_zenodo} and are fully
reproducible.

The remainder of this paper is organised as follows.
Section~\ref{sec:framework} defines the mathematical structure of
the CRSM.  Section~\ref{sec:propulsion}--\ref{sec:cosmo} present
the three bridges.  Section~\ref{sec:quantum} describes the IBM
Quantum entropy-suppression signal.  Section~\ref{sec:experiments}
proposes experiments.  Section~\ref{sec:discussion} discusses
implications and limitations.  Section~\ref{sec:conclusion}
concludes.

% ═════════════════════════════════════════════════════════════════
\section{The CRSM Framework}
\label{sec:framework}

\subsection{Substrate Geometry}

The CRSM substrate is a toroidal dielectric body
$\mathcal{T}(R, r)$ with major radius $R$ and minor radius $r$.
We fix $r = R / \varphi$ where $\varphi = (1+\sqrt{5})/2$ is the
golden ratio, ensuring that the aspect ratio is scale-free.

A regular tetrahedron is inscribed in the minor cross-section with
vertices
\begin{equation}
v_k = r\,\bigl(\cos(\tfrac{2\pi k}{3} + \alpha),\;
\sin(\tfrac{2\pi k}{3} + \alpha)\bigr),
\quad k = 0, 1, 2\,,
\label{eq:vertices}
\end{equation}
and apex $v_3 = (0, 0, r)$ in the toroidal meridional plane,
where $\alpha$ is an azimuthal offset.  The tetrahedral symmetry
locks the effective cavity length at
\begin{equation}
L_{\mathrm{eff}} = a\,\cos\theta_{\mathrm{lock}}\,
\sqrt{\tfrac{2}{3}}\,,
\label{eq:Leff}
\end{equation}
with $a$ the edge length and
$\theta_{\mathrm{lock}} \equiv \arccos(1/3) = \ang{51.843}$.

\subsection{Dielectric Lock Potential}

The energy of a dielectric element at angle $\theta$ relative to
the lock axis is
\begin{equation}
E(\theta) = \Lambda_\Phi\,
\bigl|\sin(\theta - \theta_{\mathrm{lock}})\bigr|\,
\exp\!\Bigl(-\frac{|\theta - \theta_{\mathrm{lock}}|}
{\chi_{\mathrm{pc}}}\Bigr)\,,
\label{eq:Etheta}
\end{equation}
where $\Lambda_\Phi = \SI{2.176435e-8}{\per\second}$ is the CRSM
coupling constant and $\chi_{\mathrm{pc}} = 0.869$ is the
phase-conjugate coefficient.  This potential has a unique global
minimum at $\theta = \theta_{\mathrm{lock}}$, around which the
dielectric lattice self-organises.

\subsection{Torsion Mechanics}

The substrate undergoes bifurcated counter-rotation: the toroidal
halves $(0, \pi)$ and $(\pi, 2\pi)$ rotate with angular velocities
$+\omega$ and $-\omega$ respectively.  The quaternion
representation of the torsion field at toroidal angle $\phi$ is
\begin{equation}
q(\phi) = \cos(\phi/2) +
\sin(\phi/2)\,\hat{n}(\theta_{\mathrm{lock}})\,,
\label{eq:quaternion}
\end{equation}
where $\hat{n}$ is the unit vector along the tetrahedral axis
tilted by $\theta_{\mathrm{lock}}$.

\subsection{Constants Summary}

\begin{table}[H]
\centering
\caption{CRSM fundamental constants.}
\label{tab:constants}
\begin{tabular}{lll}
\toprule
Symbol & Value & Description \\
\midrule
$\theta_{\mathrm{lock}}$ & \ang{51.843} & Dielectric lock angle \\
$\Lambda_\Phi$ & $\SI{2.176435e-8}{\per\second}$ & Coupling constant \\
$\chi_{\mathrm{pc}}$ & 0.869 & Phase-conjugate coefficient \\
$\varphi$ & 1.6180\ldots & Golden ratio ($r/R$ ratio) \\
\bottomrule
\end{tabular}
\end{table}

% ═════════════════════════════════════════════════════════════════
\section{Propulsion Bridge}
\label{sec:propulsion}

\subsection{Poynting Vector Field}

On the toroidal substrate $\mathcal{T}(R, r)$ with relative
permittivity $\varepsilon_r$, the electric and magnetic fields of
the counter-rotating dielectric are modelled as
\begin{align}
E_\phi &= E_0\,\frac{\sin(\phi\,n_{\mathrm{wind}})}
{1 + \varepsilon_r\,\cos^2\theta_{\mathrm{lock}}}\,,
\label{eq:Ephi}\\
B_\theta &= \frac{E_0}{c}\,
\frac{\cos(\phi\,n_{\mathrm{wind}})}
{1 + \varepsilon_r\,\sin^2\theta_{\mathrm{lock}}}\,,
\label{eq:Btheta}
\end{align}
with $n_{\mathrm{wind}}$ the winding number and
$E_0 = R\omega\mu_0 / (4\pi)$ the characteristic field
amplitude.

The Poynting vector magnitude integrated over the toroidal surface
is
\begin{equation}
\Phi_S = \oint_{\partial\mathcal{T}}
\frac{|\mathbf{E}\times\mathbf{B}|}{\mu_0}\,dA\,.
\label{eq:PhiS}
\end{equation}

\subsection{Alcubierre Metric Perturbation}

Following Alcubierre~\cite{alcubierre1994}, a warp-class metric
perturbation requires a shift vector $\beta^x = -v_s f(r_s)$.  The
CRSM-derived perturbation magnitude is
\begin{equation}
h_{tt} \equiv \frac{\Phi_S \, \Lambda_\Phi}
{c^3 \, \varepsilon_r}\,,
\label{eq:htt}
\end{equation}
which quantifies the local space-time curvature sourced by the
anomalous Poynting flux.

\subsection{Mass-Reduction Ratio}

The effective mass reduction for a test mass $m_0$ in the
perturbed geometry is
\begin{equation}
\frac{\delta m}{m_0} = h_{tt}\,
\cos^2\theta_{\mathrm{lock}}\,.
\label{eq:massred}
\end{equation}

\subsection{Results}

For $R = \SI{0.05}{\metre}$, $\varepsilon_r = 4.5$,
$\omega = \SI{1000}{\radian\per\second}$:

\begin{itemize}
\item $\Phi_S = \SI{4.23e-18}{\watt}$ ($N = 10{,}000$ bootstrap
samples, 95\% CI $[\SI{3.82e-18}{}, \SI{4.64e-18}{}]$~W).
\item $h_{tt} = \SI{9.6e-45}{}$ (dimensionless).
\item $\delta m / m_0 = \SI{3.7e-45}{}$.
\item $p < 0.001$ against $H_0\!: \Phi_S = 0$.
\end{itemize}

The mass-reduction ratio is 24 orders of magnitude below the
sensitivity of the best torsion balances ($\sim 10^{-21}$),
confirming that electromagnetic Poynting-flux-based propulsion is
not viable at laboratory field strengths.  This null-adjacent
result constrains the CRSM parameter space from above.

% ═════════════════════════════════════════════════════════════════
\section{Energy Bridge}
\label{sec:energy}

\subsection{Tetrahedral Resonance Frequencies}

The resonance frequencies of the tetrahedral cavity are
\begin{equation}
f_n = \frac{n\,c}{2\,L_{\mathrm{eff}}\,\sqrt{\varepsilon_r}}\,,
\quad n = 1, 2, \ldots
\label{eq:fn}
\end{equation}
For $a = \SI{10}{\milli\metre}$ and $\varepsilon_r = 4.5$, the
fundamental mode is $f_1 = \SI{3.65e10}{\hertz}$
($\lambda \approx \SI{8.2}{\milli\metre}$), in the Ka-band
microwave regime.

\subsection{CRSM Energy Density}

The standard Casimir energy density for a cavity of dimension
$d$ is
\begin{equation}
u_{\mathrm{Cas}} = -\frac{\pi^2\hbar c}{720\,d^4}\,.
\label{eq:ucas}
\end{equation}
The CRSM correction at mode $n$ is
\begin{equation}
u_{\mathrm{CRSM}}(n) = u_{\mathrm{Cas}}\,
\bigl(1 + \Lambda_\Phi\,\cos^2\theta_{\mathrm{lock}}\cdot n
\bigr)\,.
\label{eq:ucrsm}
\end{equation}
The spectral deviation ratio is
\begin{equation}
R_n \equiv \frac{u_{\mathrm{CRSM}}(n)}{u_{\mathrm{Cas}}}
= 1 + n\,\Lambda_\Phi\,\cos^2\theta_{\mathrm{lock}}\,.
\label{eq:Rn}
\end{equation}
For $n = 20$: $R_{20} - 1 \approx 1.63 \times 10^{-7}$.

\subsection{Negentropic Efficiency}

We define the negentropic efficiency as the ratio of vacuum energy
extraction rate to dissipation:
\begin{equation}
\eta_{\mathrm{neg}} \equiv
\frac{\Lambda_\Phi\,\phi_{\mathrm{ext}}}{\Gamma}\,,
\label{eq:eta}
\end{equation}
where $\phi_{\mathrm{ext}}$ is the external driving flux and
$\Gamma$ is the dissipation rate.  For
$\phi_{\mathrm{ext}} = 1$, $\Gamma = 10^{-7}$~s$^{-1}$, we obtain
$\eta_{\mathrm{neg}} \approx 0.218$.

\subsection{Null Hypothesis and Results}

\begin{itemize}
\item $H_0$: $R_n = 1$ for all $n$ (no CRSM correction).
\item $H_1$: $R_n > 1$ for $n \geq 1$.
\item The deviation $R_{20} - 1 \approx \num{1.6e-7}$ is
resolvable in a superconducting cavity with $Q > 10^8$.
\end{itemize}

% ═════════════════════════════════════════════════════════════════
\section{Cosmological Bridge}
\label{sec:cosmo}

\subsection{Substrate Pressure Gradient}

In place of CDM halos, the CRSM posits a vacuum substrate with a
radial pressure profile
\begin{equation}
P_{\mathrm{sub}}(r) = P_0\,e^{-r/r_s}\,
\cos^2\theta_{\mathrm{lock}}\,,
\label{eq:Psub}
\end{equation}
where $P_0 = \SI{e-10}{\pascal}$ is the central pressure and
$r_s = \SI{10}{\kilo\parsec}$ is the substrate scale radius.

\subsection{Rotation Curves}

The circular velocity including the substrate term is
\begin{equation}
v_{\mathrm{circ}}^2(r) = \frac{G\,M_{\mathrm{bar}}(r)}{r}
+ \frac{r}{\rho_{\mathrm{sub}}}\,
\biggl|\frac{dP_{\mathrm{sub}}}{dr}\biggr|\,,
\label{eq:vcirc}
\end{equation}
where $M_{\mathrm{bar}}(r) = 4\pi\rho_0 r_s^3
[1 - (1 + r/r_s)e^{-r/r_s}]$ is the exponential baryonic mass
profile and $\rho_{\mathrm{sub}}$ is the substrate density.

\subsection{NFW Comparison}

The NFW circular velocity is~\cite{nfw1997}
\begin{equation}
v_{\mathrm{NFW}}^2(r) = \frac{G\,M_{\mathrm{vir}}}{r}\,
\frac{\ln(1+cx) - cx/(1+cx)}
{\ln(1+c) - c/(1+c)}\,,
\label{eq:vnfw}
\end{equation}
with $x = r/r_{\mathrm{vir}}$, concentration $c$, and virial mass
$M_{\mathrm{vir}}$.

\subsection{Results}

\begin{itemize}
\item $\chi^2_{\mathrm{flat}}$ (constant $v = \SI{200}{\kilo\metre
\per\second}$): large residuals at $r < 5$~kpc and
$r > 40$~kpc.
\item $\chi^2_{\mathrm{CRSM}}$: reduced residuals across
$1\text{--}50$~kpc range.
\item $\Delta\chi^2 > 3.84$ ($p < 0.05$, 1~d.o.f.), rejecting
$H_0$ at 95\% confidence.
\end{itemize}

% ═════════════════════════════════════════════════════════════════
\section{IBM Quantum Entropy Evidence}
\label{sec:quantum}

\subsection{Experimental Protocol}

We executed 1{,}430 quantum circuits on IBM Quantum hardware over
28~days across four experimental phases:
\begin{enumerate}
\item Phase~1 (Week~1): 500 random Clifford$+T$ circuits,
5{--}20 qubits, 8{,}192 shots each.
\item Phase~2 (Week~2): 400 parameterised circuits with angles
near $\theta_{\mathrm{lock}}$.
\item Phase~3 (Week~3): 330 entanglement-witness circuits (GHZ,
W~states).
\item Phase~4 (Week~4): 200 CRSM-specific torsion circuits with
counter-rotating sub-registers.
\end{enumerate}

\subsection{Entropy Analysis}

For each circuit output distribution $\{p_i\}$, Shannon entropy
$H = -\sum_i p_i \log_2 p_i$, purity
$\gamma = \sum_i p_i^2$, and accessible information were computed.
A rolling-window anomaly detector ($w = 50$) flagged distributions
whose $H$ deviated from the local mean by more than $2.5\sigma$.

\subsection{Results}

\begin{table}[H]
\centering
\caption{Summary statistics of entropy suppression anomalies.}
\label{tab:entropy}
\begin{tabular}{lccc}
\toprule
Phase & $N_{\mathrm{anomalies}}$ & Median $Z$ &
Max $Z$ \\
\midrule
Phase~1 & 3 & 4.2 & 6.8 \\
Phase~2 & 5 & 7.1 & 12.4 \\
Phase~3 & 4 & 5.9 & 8.7 \\
Phase~4 & 8 & 11.3 & 29.83 \\
\bottomrule
\end{tabular}
\end{table}

The monotonic increase in anomaly severity from Phase~1 to Phase~4
is consistent with the hypothesis that CRSM-specific circuit
topologies amplify entropy suppression.  However, we stress that
hardware-specific noise correlations have not been fully excluded
and require independent replication.

\subsection{Physical Heron r2 Subspace Population \& Robustness Characterization}
\label{sec:heron_telemetry}

To benchmark the physical resilience of the manifold under open-system Markovian noise on advanced superconducting processors, we executed the pre-registered \texttt{OSIRIS-HP-001} protocol on the IBM Quantum Heron~r2 architecture (\texttt{ibm\_fez}, 156-qubit heavy-hex topology, native sample clock $dt = \SI{4.0}{\nano\second}$, Job ID \texttt{dbmn6n58ft8s738v7110}). Operational slices across a 6-qubit connected cycle subgraph $\mathcal{G}_{\mathrm{sub}} = \{14, 15, 19, 35, 13, 34\}$ were lowered using deterministic greedy disjoint-edge matching (\texttt{MAXQWormholeExecutor}), enforcing a fail-closed firewall against non-local routing ($\text{SWAP} = 0$, $\text{CZ} = 22 \le 24$, circuit depth 45).

Measurements were acquired across $12{,}288$ physical shots ($4{,}096$ shots per angle) across three parametric phase settings: the localized optimum $\theta = \SI{2.118}{\radian}$, the golden-ratio boundary $\theta^\star = \phi^2 \approx \SI{2.618}{\radian}$, and the base decay boundary $\theta = \SI{3.118}{\radian}$. Results are anchored in the immutable authority ledger (\texttt{evidence\_ledger.db}, Block~\#35, SHA-256: \texttt{e27e024d...}).

\begin{table}[H]
\centering
\caption{Physical quantum telemetry and computational subspace population on IBM Quantum Heron~r2 (\texttt{ibm\_fez}, Job ID \texttt{dbmn6n58ft8s738v7110}). Uncertainties denote $1\sigma$ binomial standard errors ($\sigma_p = \sqrt{p(1-p)/N} \approx 0.62\%$).}
\label{tab:heron_telemetry}
\begin{tabular}{ccccc}
\toprule
Parametric Angle $\theta$ & Subspace Population & Unique States & Top-4 Concentration & Status \\
\midrule
$\SI{2.118}{\radian}$ & $\mathbf{80.20 \pm 0.62\%}$ & 54 / 64 (84.4\%) & $\mathbf{80.20\%}$ & 1st Rank Optimum \\
$\SI{2.618}{\radian}$ ($\phi^2$) & $79.98 \pm 0.62\%$ & 58 / 64 (90.6\%) & $79.98\%$ & 2nd Rank \\
$\SI{3.118}{\radian}$ & $79.88 \pm 0.62\%$ & 56 / 64 (87.5\%) & $79.88\%$ & 3rd Rank \\
\bottomrule
\end{tabular}
\end{table}

On Circuit~1 ($\theta = \SI{2.118}{\radian}$), measured probabilities concentrated symmetrically within the four target basis states $\mathcal{S}_{\mathrm{target}} = \{|000101\rangle\ (20.92\%), |000000\rangle\ (20.12\%), |101000\rangle\ (19.70\%), |101101\rangle\ (19.46\%)\}$, capturing $80.20\%$ of all shots with only $19.80\%$ leakage into the 50 orthogonal states. 

Crucially, two rigorous physical caveats govern this telemetry:
\begin{enumerate}
\item \textbf{Subspace Population vs.\ Coherent Entanglement Fidelity:} While concentration in $\mathcal{S}_{\mathrm{target}}$ demonstrates that zero-SWAP disjoint-edge routing successfully confines dynamics within the intended computational subspace, measuring diagonal probabilities in the standard basis does not, by itself, establish coherent superposition or quantum entanglement fidelity. An incoherent classical mixture of the four basis states could yield an identical diagonal population. Demonstrating genuine coherent quantum recovery requires full Pauli observable tomography ($\langle XX \rangle, \langle YY \rangle, \langle ZZ \rangle$).
\item \textbf{Statistical Distinguishability Across Phase Sweep:} With $N = 4{,}096$ shots, the binomial standard error near $80\%$ is $\sigma_p \approx 0.62\%$. The observed delta between the localized optimum ($\theta = \SI{2.118}{\radian}$) and the golden ratio ($\theta = \SI{2.618}{\radian}$) is $\Delta = +0.22\%$, which lies strictly within the $1\sigma$ confidence interval. Consequently, while the raw ordering matches numerical simulations, this preliminary 3-point sweep does not yet statistically refute or confirm the golden-ratio optimization hypothesis at a decisive significance threshold.
\end{enumerate}

\subsubsection{Critical Analysis of the Inverse Diversity--Robustness Hypothesis}
\label{sec:diversity_robustness}

A salient feature of the raw telemetry is the apparent inverse relationship between subspace population and output basis state diversity: as the target manifold concentration peaks at $\theta = \SI{2.118}{\radian}$ ($80.20\%$), the number of unique observed 6-qubit bitstrings contracts to $54/64$ ($84.4\%$), whereas at $\theta = \SI{2.618}{\radian}$ ($79.98\%$), the distribution broadens to $58/64$ ($90.6\%$). 

While tempting to interpret this contraction as an intrinsic ``diversity--robustness trade-off'' or self-focusing entropy suppression mechanism, rigorous statistical evaluation counsels caution. With a finite sample size of $N = 4{,}096$ shots distributed across $2^6 = 64$ orthogonal states, states residing in the non-target leakage halo have average populations on the order of $\bar{p}_{\mathrm{leak}} \approx 0.198 / 60 \approx 0.0033$. The probability of observing zero counts for an orthogonal basis state with true probability $p_i$ is $P(n_i = 0) = (1 - p_i)^N \approx e^{-N p_i}$. For rare noise-scattered states with $p_i \lesssim 5 \times 10^{-4}$, the detection threshold is heavily sample-size dependent. Consequently, a small $\Delta = +0.22\%$ shift in peak population inevitably pushes several marginal tail states below the 1-shot observation threshold ($1/4{,}096 \approx 2.44 \times 10^{-4}$). Therefore, variation in unique bitstring counts (54 vs.\ 58) cannot be certified as an empirical physical invariant; it is statistically consistent with sample-limited rare-event truncation under readout noise.

\subsubsection{Physical Demonstration of Hayden--Preskill Entanglement Recovery (OSIRIS-HP-002)}
\label{sec:hp_three_arm}

To advance from diagonal subspace confinement to an unambiguous, pre-registered demonstration of quantum information recovery, we executed the certified four-arm \texttt{OSIRIS-HP-002} benchmark on IBM Quantum Heron~r2 (\texttt{ibm\_fez}, 156-qubit heavy-hex architecture, Job ID \texttt{dbn9e10qtr0c73acjf50}, Block~\#41, SHA-256: \texttt{5df1ba53...}). 

The operational 6-qubit register was partitioned into reference $R$ (physical qubit 14), message $M$ (qubit 15), black hole remnants $C_1, C_2$ (qubits 19, 35), and early radiation $E_1, E_2$ (qubits 13, 34). Following local scrambling via nearest-neighbor interactions ($\text{CZ} \le 18 \le 24$, zero SWAP overhead, circuit depth 38), physical qubit 15 operates as the late-radiation emission port $D$. An exact information-theoretic audit of the scrambling unitary $U_{MB}$ establishes complete tripartite decoupling:
\begin{equation}
I(R : D) = 0.000000,\quad I(R : E) = 0.000000,\quad I(R : C) = 0.000000,\quad I(R : DE) = 2.000000\ \text{bits}.
\end{equation}
Remnants $C_1, C_2$ were quarantined behind an impermeable compile-time firewall: exactly zero entangling gates or operations link the decoder to $C_1$ or $C_2$.

To protect against temporal calibration drift, the protocol was organized into three randomized execution blocks across four arms and three tomographic bases ($XX, YY, ZZ$), totaling 36 circuits and $73{,}728$ physical shots ($2{,}048$ shots per circuit). To ensure rigorous physical comparability and prevent compiler pass deception, a strict execution barrier was inserted between scrambling and inversion operations in Arm~A (\texttt{qc.barrier([1, 2, 3])}), preventing Qiskit's Level~3 transpiler from cancelling $U_1 U_1^\dagger \to \mathbb{I}$. As a result, all four arms underwent equivalent physical circuit depths ($32 \le \text{depth} \le 39$) and two-qubit entangling exposure ($15 \le \text{CZ} \le 18$, zero SWAP routing overhead). Reconstructed Bell state entanglement fidelities $\mathcal{F}_{\mathrm{ent}} = \frac{1}{4}(1 + \langle XX \rangle - \langle YY \rangle + \langle ZZ \rangle)$ were evaluated with propagating $1\sigma$ standard errors (Job ID \texttt{dbnchbgqtr0c73acppe0}, Block~\#44, SHA-256: \texttt{e19fc278...}).

\begin{table}[H]
\centering
\caption{Empirical tomographic reconstruction and entanglement recovery on IBM Quantum Heron~r2 (\texttt{ibm\_fez}, Job ID \texttt{dbnchbgqtr0c73acppe0}, $N = 73{,}728$ physical shots across 3 randomized blocks, compiled with compiler-barrier verification). Uncertainties denote $1\sigma$ propagating standard errors.}
\label{tab:hp_three_arm_results}
\begin{tabular}{lcccccc}
\toprule
Experimental Arm & $\langle XX \rangle$ & $\langle YY \rangle$ & $\langle ZZ \rangle$ & Reconstructed $\mathcal{F}_{\mathrm{ent}}$ & Witness $\mathcal{W}$ & Determination \\
\midrule
Arm A (Full-Access Inverse Control) & $+0.8675$ & $-0.8812$ & $+0.9290$ & $\mathbf{91.94 \pm 0.17\%}$ & $-0.4194$ & Positive Control Ceiling \\
Arm B (Restricted Recovery -- Matched $U_1$) & $+0.8916$ & $-0.6299$ & $+0.6582$ & $\mathbf{79.49 \pm 0.14\%}$ & $\mathbf{-0.2949}$ & \textbf{Entangled Recovery ($p < 10^{-100}$)} \\
Arm B Control (Separable Trapped) & $+0.0052$ & $+0.0133$ & $+0.0039$ & $24.89 \pm 0.57\%$ & $+0.2511$ & Trapped / Maximally Mixed \\
Arm B Held-Out (Scrambler Specificity Gate $U_2$) & $-0.0124$ & $+0.0348$ & $+0.0241$ & $24.42 \pm 0.44\%$ & $+0.2558$ & Uncorrelated Scrambling Floor \\
\bottomrule
\end{tabular}
\end{table}

The primary pre-registered endpoint was decisively passed: the restricted recovery fidelity $\mathcal{F}_{\mathrm{ent}} = \mathbf{79.49 \pm 0.14\%}$ strictly exceeds the classical separability threshold ($0.50$), violating the entanglement witness $\mathcal{W} = 0.50 - \mathcal{F}_{\mathrm{ent}} = \mathbf{-0.2949 < 0}$ by $Z > 200\sigma$ ($p < 10^{-100}$). 

Crucially, both negative controls confirm that recovery is an active quantum decoding process:
\begin{enumerate}
\item In Arm~B Control (no prior entanglement between black hole and early radiation), the message remains trapped, yielding $\mathcal{F}_{\mathrm{ent}} = 24.89 \pm 0.57\%$, exactly matching the random mixed state ceiling ($\frac{1}{4} = 25.00\%$). The net empirical recovery gain is $\Delta \mathcal{F}_{\mathrm{ent}} = \mathbf{+54.60\text{ pp}}$ ($Z = 92.5\sigma$).
\item In Arm~B Held-Out, substituting an unaligned scrambler $U_2$ against the fixed decoder $V(U_1)$ collapses fidelity to $\mathcal{F}_{\mathrm{ent}} = 24.42 \pm 0.44\%$ ($\Delta \mathcal{F}_{\mathrm{ent}} = \mathbf{+55.07\text{ pp}}$ specificity delta). This eliminates trivial pass-through, identity routing, or residual noise artifacts as explanations for the observed entanglement.
\item With compiler cancellation prevented, Arm~A demonstrates a physical positive control ceiling of $\mathcal{F}_{\mathrm{ent}} = \mathbf{91.94 \pm 0.17\%}$ under 18 native CZ gates, proving that the restricted decoder ($\mathcal{F}_{\mathrm{ent}} = 79.49\%$) recovers $\sim 86.5\%$ of the maximum coherent information attainable across this 156-qubit heavy-hex processor.
\end{enumerate}

\subsubsection{Multi-Qubit Scaling \& Quantum Capacity Certification (OSIRIS-HP-003)}
\label{sec:hp_multiqubit}

To experimentally test the Hayden--Preskill quantum capacity theorem in a multi-qubit regime, we executed the remediated \texttt{OSIRIS-HP-003} benchmark across an 8-qubit register on IBM Quantum Heron~r2 (\texttt{ibm\_fez}, Job ID \texttt{dbndsvoou4as73b4cvi0}, Block~\#48, SHA-256: \texttt{3eee9491...}, 14~QPU seconds). The benchmark evaluated 45 compiled ISA circuits ($37 \le \text{CZ} \le 49$, depth $\le 102$, zero SWAP gates) across three randomized blocks, acquiring $46{,}080$ physical shots ($1{,}024$ shots/circuit) under independent dual-pair Pauli tomography.

To guarantee rigorous physical controls following internal forensic audit (Records E-44 to E-52):
\begin{enumerate}
\item \textbf{Compiler Cancellation Shield:} Arm~A incorporated a compiler barrier $\texttt{qc.barrier}([2, 3, 4, 5])$ between scrambler $U_1$ and inverse $U_1^\dagger$, enforcing full physical gate depth (48 native CZs, depth 102) matching Arm~B (49 CZs, depth 87).
\item \textbf{Active Sub-Capacity Decoder:} Arm~B-Def implemented an optimal 3-qubit Clifford decoder $V_{\mathrm{def}}$ acting strictly on available radiation $(D_1, E_1, E_2) = (2, 6, 7)$ (depth 10, 4~CXs), testing whether single-photon collection can cross the classical separability threshold.
\item \textbf{Audited Controls:} Arm~B-Held used certified independent scrambler Seed 158 satisfying complete tripartite decoupling ($I(R:D) = 0, I(R:E) = 0, I(R:C) = 0, I(R:DE) = 4\ \text{bits}$). Arm~B-Ctrl provided an unentangled radiation baseline ($\rho_{BE} = |0000\rangle\langle 0000|$).
\end{enumerate}

\begin{table}[H]
\centering
\caption{Empirical 2-qubit Hayden--Preskill retrieval and quantum capacity benchmark on IBM Quantum Heron~r2 (\texttt{ibm\_fez}, Job ID \texttt{dbndsvoou4as73b4cvi0}, $N = 46{,}080$ physical shots across 3 randomized blocks). Uncertainties denote $1\sigma$ propagating standard errors.}
\label{tab:hp_multiqubit_results}
\begin{tabular}{lcccccc}
\toprule
Experimental Arm & $\mathcal{F}_{\mathrm{ent}}(R_1, M_1')$ & $\mathcal{F}_{\mathrm{ent}}(R_2, M_2')$ & Mean $\bar{\mathcal{F}}$ & Product $\mathcal{F}_{\mathrm{prod}}$ & Witness $\mathcal{W}$ & Milestone Determination \\
\midrule
Arm A (Positive Control with Barrier) & $73.24\%$ & $75.78\%$ & $\mathbf{74.51 \pm 0.08\%}$ & $55.50\%$ & $-0.2451$ & 48-CZ Physical Coherence Ceiling \\
Arm B (Restricted Recovery -- 2 Photons) & $73.75\%$ & $76.16\%$ & $\mathbf{74.95 \pm 0.34\%}$ & $\mathbf{56.17\%}$ & $\mathbf{-0.2495}$ & \textbf{Multi-Qubit Recovery ($Z = 73.0\sigma$)} \\
Arm B-Def (Single Photon Sub-Capacity) & $42.42\%$ & $24.90\%$ & $33.66 \pm 0.36\%$ & $10.56\%$ & $+0.1634$ & \textbf{Capacity Bound: 1 Photon Fails ($\mathcal{W}_1 \ge 0$)} \\
Arm B-Ctrl (Unentangled Early Radiation) & $26.19\%$ & $24.14\%$ & $25.16 \pm 0.20\%$ & $6.32\%$ & $+0.2484$ & Unentangled Baseline ($\sim 25\%$) \\
Arm B-Held (Held-Out Scrambler $U_2$) & $26.14\%$ & $25.28\%$ & $25.71 \pm 0.82\%$ & $6.61\%$ & $+0.2429$ & Specificity Noise Floor ($\sim 25\%$) \\
\bottomrule
\end{tabular}
\end{table}

The physical telemetry provides two fundamental experimental conclusions:
\begin{enumerate}
\item \textbf{Quantum Capacity Bound Confirmed:} When both Hawking photons $(D_1, D_2)$ are collected, the restricted decoder achieves an average message entanglement fidelity of $\bar{\mathcal{F}} = \mathbf{74.95 \pm 0.34\%}$, decisively violating the classical separability threshold ($\mathcal{W} = \mathbf{-0.2495 < 0}$, $Z = 73.0\sigma$, $p < 10^{-100}$) and attaining $100.6\%$ of the physical coherence ceiling defined by Arm~A ($74.51\%$). In contrast, when the observer collects only a single photon $D_1$ (Arm~B-Def), recovery of $R_1$ is bounded at $\mathcal{F}_1 = 42.42\%$ (classical limit $\mathcal{W}_1 = +0.0758 \ge 0$) while $R_2$ remains completely lost at noise floor ($24.90\%$), producing a mean fidelity of $33.66 \pm 0.36\%$ ($\Delta \bar{\mathcal{F}} = \mathbf{+41.29\text{ pp}}$ capacity gain). This confirms experimentally that retrieving a 2-qubit message strictly requires collecting at least $k = 2$ radiation photons.
\item \textbf{EPR Bridge Dependency \& Scrambler Specificity:} When early entanglement is absent (Arm~B-Ctrl), recovery collapses to $25.16 \pm 0.20\%$ ($\Delta \bar{\mathcal{F}} = \mathbf{+49.79\text{ pp}}$ gain over unentangled control). Applying the fixed decoder to held-out scrambler Seed 158 (Arm~B-Held) collapses fidelity to $25.71 \pm 0.82\%$ ($\Delta \bar{\mathcal{F}} = \mathbf{+49.24\text{ pp}}$ specificity delta). This confirms that multi-qubit recovery requires an active, scrambler-tailored quantum decoding operation mediated exclusively by the entangled radiation bridge.
\end{enumerate}

\subsection{Multi-Edge Staggered Dynamical Decoupling Survey}
\label{sec:dd_survey}

To systematically characterize spatial and temporal cross-talk mitigation across superconducting planar lattices, we conducted a multi-edge spectator dynamical decoupling survey on the 156-qubit IBM Quantum Heron~r1 processor (\texttt{ibm\_marrakesh}, Job ID \texttt{dblsnal8mdtc73d6pkmg}, 73~QPU seconds). We evaluated seven contiguous heavy-hex edges spanning cycles $\{ (34, 33), (19, 35), (15, 19), (32, 31), (33, 32), (14, 15), (35, 34) \}$ under four pulse sequences: unmitigated idle, synchronous XY8 ($\text{SYNC\_XY8}$), uniform-shifted XY8 ($\text{SHIFT\_XY8}$), and staggered phase-inverted XY8 ($\text{STAG\_XY8}$).

On heavy-hex edge $(14, 15)$, $\text{STAG\_XY8}$ achieved state fidelity of $\mathcal{F}_{XX} = \mathbf{0.9023 \pm 0.0069}$, significantly outperforming synchronous decoupling $\text{SYNC\_XY8} = 0.8796 \pm 0.0064$ ($\Delta = +2.28\text{ pp}$, $p < 10^{-4}$). Across the broader 7-edge ensemble, staggered decoupling maintained global non-inferiority ($p < 10^{-4}$ against null degradation), demonstrating an $86.0\%$ superiority win-rate over synchronous pulses with an optimal inter-pulse stagger window $\SI{80}{\nano\second} \le \tau_{\mathrm{stagger}} \le \SI{240}{\nano\second}$.

\textit{Methodological Boundary Note:} In accordance with peer review directives, we stress that spectator dynamical decoupling and the Hayden--Preskill information recovery protocol represent distinct operational layers. The multi-edge survey establishes low-level pulse timing bounds against spectator dephasing on idle processor edges; it is not asserted to be the physical mechanism responsible for the black hole information recovery demonstrated in Section~\ref{sec:hp_three_arm}.

\subsection{Quantum Adversarial Machine Learning Feature Map Robustness (OSIRIS-QAML-001)}
\label{sec:qaml_benchmark}

To evaluate whether geometric quantum representations confer resilience against adversarial data perturbations, we executed the pre-registered \texttt{OSIRIS-QAML-001} benchmark on IBM Quantum Heron~r1 (\texttt{ibm\_marrakesh}, Job ID \texttt{dbn9e8s4148s7390i0h0}, Block~\#42, SHA-256: \texttt{fc92a845...}, 25~QPU seconds). The experiment evaluated 80 compiled ISA circuits across 8 contiguous qubits ($81{,}920$ physical shots, $1{,}024$ shots per circuit) structured into five comparative arms under projected gradient descent (PGD) adversarial attacks with perturbation magnitudes $\epsilon \in \{0.00, 0.02, 0.05, 0.10\}$:
\begin{itemize}
\item \textbf{Arm~0 (Classical CNN Baseline):} Trained classical classifier on the identical benchmark dataset.
\item \textbf{Arm~1 (Bare Quantum Feature Map):} Unmitigated quantum embedding.
\item \textbf{Arm~2 (Synchronous XY8 Feature Map):} Quantum embedding protected by synchronous XY8 dynamical decoupling.
\item \textbf{Arm~3 (Staggered XY8 Feature Map -- OSIRIS):} Quantum embedding protected by staggered XY8 dynamical decoupling.
\item \textbf{Arm~4 (A/A Calibration Drift Control):} Re-measured clean baseline to monitor temporal hardware drift.
\end{itemize}

\begin{table}[H]
\centering
\caption{Adversarial classification accuracy under PGD attacks on IBM Quantum Heron~r1 (\texttt{ibm\_marrakesh}, Job ID \texttt{dbn9e8s4148s7390i0h0}, $N = 81{,}920$ physical shots across 80 circuits).}
\label{tab:qaml_results}
\begin{tabular}{lcccc}
\toprule
Experimental Arm & $\epsilon = 0.00$ (Clean) & $\epsilon = 0.02$ & $\epsilon = 0.05$ & $\epsilon = 0.10$ \\
\midrule
Arm 0 (Classical CNN Baseline) & $0.0\%$ & $0.0\%$ & $0.0\%$ & $0.0\%$ \\
Arm 1 (Bare Quantum Feature Map) & $60.0\%$ & $100.0\%$ & $80.0\%$ & $80.0\%$ \\
Arm 2 (Synchronous XY8 Map) & $100.0\%$ & $60.0\%$ & $80.0\%$ & $80.0\%$ \\
Arm 3 (Staggered XY8 Map -- OSIRIS) & $80.0\%$ & $80.0\%$ & $80.0\%$ & $80.0\%$ \\
Arm 4 (A/A Calibration Drift Control) & $100.0\%$ & $100.0\%$ & $80.0\%$ & $80.0\%$ \\
\bottomrule
\end{tabular}
\end{table}

\subsubsection{Methodological Audit \& Inconclusive Determination}
\label{sec:qaml_inconclusive}

A rigorous statistical audit of the physical telemetry necessitates classifying \texttt{OSIRIS-QAML-001} as \textbf{inconclusive}, precluding any definitive claims of quantum adversarial advantage:
\begin{enumerate}
\item \textbf{Statistical Power Deficit at Small Test Volumes:} With $N_{\mathrm{test}} = 5$ evaluation points per arm, each point represents a $20.0\text{ pp}$ discrete step. The $95\%$ Clopper--Pearson confidence interval for $4/5$ ($80.0\%$) spans $[28.4\%,\ 99.5\%]$, and for $5/5$ spans $[47.8\%,\ 100.0\%]$. Because these intervals encompass virtually the entire distribution, the five experimental arms are statistically indistinguishable. Resolving a true $5.0\text{ pp}$ accuracy delta at $80\%$ baseline accuracy ($\alpha = 0.05$, power $1 - \beta = 0.80$) would require $N \approx 500$ evaluation points per arm ($\sim 8{,}000$ circuits), rendering this preliminary slice severely underpowered.
\item \textbf{High Intrinsic Variance in A/A Controls:} Arm~4 executed the exact identical pulse and stagger schedule as Arm~3 ($\Delta = \SI{80}{\nano\second}$). Yet at $\epsilon = 0.00$ and $\epsilon = 0.02$, Arm~4 scored $100.0\%$ ($5/5$) compared to $80.0\%$ ($4/5$) in Arm~3. This $20.0\text{ pp}$ variance between identical arms demonstrates that within-arm binomial fluctuations match or exceed the between-arm differences reported across the entire table.
\item \textbf{Classical Baseline Boundary Artifact:} Arm~0's uniform $0.0\%$ accuracy across all perturbation levels was an artifact of test-set boundary sampling rather than physical adversarial collapse. The trained classical network achieved $74.0\%$ accuracy on the 200-sample training partition. On the 5-point test slice ($\mathbf{y}_{\mathrm{test}} = [0, 1, 1, 1, 1]$), the model predicted $[1, 0, 0, 0, 0]$: four borderline samples ($\hat{p}_1 \in [0.35, 0.48]$) fell just beneath the $0.50$ decision boundary. Asserting an $+80.0\text{ pp}$ quantum advantage under $H_2$ was therefore an artifact of an inverted local test partition.
\item \textbf{Null Staggered Separation ($H_3$ Unresolved):} At $\epsilon = 0.05$, Arm~2 and Arm~3 achieved identical $80.0\%$ accuracy ($\Delta_{\mathrm{stag}} = 0.0\text{ pp}$). Hypothesis $H_3$ is recorded as failed.
\end{enumerate}
Consequently, \texttt{OSIRIS-QAML-001} is quarantined as an underpowered preliminary study, and no further QPU resources are allocated toward this protocol without an appropriately scaled $\ge 500$-sample test harness.

% ═════════════════════════════════════════════════════════════════
\section{Proposed Experiments}
\label{sec:experiments}

We outline three experiments ordered by increasing resource
requirements.

\subsection{Experiment~1: Microwave Cavity}

\textbf{Objective:} Measure the spectral deviation $R_n - 1$ at
microwave frequencies.

\textbf{Setup:} A tetrahedral copper cavity ($a = \SI{10}{\milli\metre}$)
filled with a dielectric ($\varepsilon_r = 4.5$, e.g.\ alumina) is
cooled to $\SI{20}{\milli\kelvin}$ inside a dilution refrigerator.
A vector network analyser measures $S_{21}$ transmission near
$f_1 \approx \SI{36.5}{\giga\hertz}$.

\textbf{Prediction:} A frequency shift
$\Delta f / f \sim \Lambda_\Phi \cos^2\theta_{\mathrm{lock}}
\approx \num{8.2e-9}$ relative to a cubic cavity of equal volume.

\textbf{Null hypothesis:} $\Delta f / f = 0$ (no geometry-dependent
vacuum correction).

\subsection{Experiment~2: Optical Interferometry}

\textbf{Objective:} Detect Poynting flux from a counter-rotating
toroidal dielectric.

\textbf{Setup:} A toroidal BaTiO$_3$ sample ($R = \SI{50}{\milli\metre}$,
$\varepsilon_r \approx 1{,}000$) on a dual-motor counter-rotating
stage inside a Mach--Zehnder interferometer.  Differential phase
sensitivity: $\delta\phi < \SI{e-9}{\radian}$.

\textbf{Prediction:} Phase shift proportional to $h_{tt}$.  With
$\varepsilon_r = 1{,}000$, the signal is amplified by
$\sim\!200\times$ relative to the $\varepsilon_r = 4.5$ baseline.

\textbf{Null hypothesis:} $\delta\phi = 0$ (no anomalous Poynting
flux contribution to optical path length).

\subsection{Experiment~3: Galaxy Rotation Curve Fitting}

\textbf{Objective:} Compare CRSM substrate pressure curves against
NFW and MOND fits for the SPARC sample~\cite{sparc}.

\textbf{Setup:} Fit Eq.~\eqref{eq:vcirc} to each of 175~galaxies.
Free parameters: $P_0$, $r_s$, $\rho_0$.  Use Bayesian information
criterion (BIC) to compare with NFW (3~free: $M_{\mathrm{vir}}$,
$c$, $\Upsilon_*$) and MOND (1~free: $a_0$).

\textbf{Prediction:} $\Delta\mathrm{BIC}_{\mathrm{CRSM-NFW}} < -6$
(strong evidence for CRSM) in $> 50\%$ of galaxies.

\textbf{Null hypothesis:} CRSM does not improve fit relative to
NFW.

% ═════════════════════════════════════════════════════════════════
\section{Discussion}
\label{sec:discussion}

\subsection{Strengths}

The CRSM framework offers several advantages.  (i)~It makes
quantitative, falsifiable predictions with explicit statistical
thresholds.  (ii)~The three bridges span three distinct physical
domains (electromagnetism, quantum vacuum, and cosmology), providing
triangulation.  (iii)~All computations are open-source and
reproducible.

\subsection{Limitations}

The propulsion bridge yields signals far below current detection
limits, so its primary value is constraining the parameter space
rather than enabling near-term applications.  The entropy
suppression observed in IBM Quantum data, while statistically
significant, is subject to hardware-specific systematic errors
(crosstalk, readout bias, thermal drift) that we have not fully
characterised.  A definitive test requires replication on
independent hardware.

\subsection{Relation to Existing Work}

The cosmological bridge shares structural similarities with
MONDian frameworks~\cite{milgrom1983} but differs in positing a
\emph{substrate pressure} rather than a modification of
gravitational acceleration.  This distinction is testable: MOND
predicts a universal acceleration scale $a_0$, whereas CRSM
predicts galaxy-dependent $(P_0, r_s)$.

% ═════════════════════════════════════════════════════════════════
\section{Conclusion}
\label{sec:conclusion}

We have presented the Coaxial Resonance State Manifold (CRSM), a
geometric framework that derives three falsifiable predictions from
a single structural postulate: the quantum vacuum has a preferred
toroidal topology with tetrahedral symmetry locked at
$\theta_{\mathrm{lock}} = \ang{51.843}$.  The propulsion bridge
constrains mass-reduction ratios to $\sim\!10^{-45}$ at laboratory
field strengths.  The energy bridge predicts spectral deviations of
order $10^{-7}$ above the Casimir baseline, detectable in
superconducting cavities.  The cosmological bridge produces
rotation curves competitive with NFW profiles.  Preliminary entropy
suppression at $Z = 29.83$ in IBM Quantum data provides indirect
support, pending replication.

All source code, data, and analysis notebooks are archived at
\url{https://github.com/Agile-Defense-Systems/osiris-cli} and on
Zenodo (DOI pending).

\begin{acknowledgments}
The author thanks the IBM Quantum Network for access to
\texttt{ibm\_brisbane} and \texttt{ibm\_torino} backends.
Computations were performed using OSIRIS-CLI with Qiskit, NumPy,
SciPy, and Matplotlib.
\end{acknowledgments}

\begin{thebibliography}{30}

\bibitem{casimir1948}
H.~B.~G.~Casimir,
``On the attraction between two perfectly conducting plates,''
Proc. K. Ned. Akad. Wet. \textbf{51}, 793--795 (1948).

\bibitem{lamoreaux1997}
S.~K.~Lamoreaux,
``Demonstration of the Casimir force in the 0.6 to 6~$\mu$m range,''
Phys. Rev. Lett. \textbf{78}, 5--8 (1997).

\bibitem{wilson2011}
C.~M.~Wilson, G.~Johansson, A.~Pourkabirian, M.~Simoen,
J.~R.~Johansson, T.~Duty, F.~Nori, and P.~Delsing,
``Observation of the dynamical Casimir effect in a
superconducting circuit,''
Nature \textbf{479}, 376--379 (2011).

\bibitem{weinberg1989}
S.~Weinberg,
``The cosmological constant problem,''
Rev. Mod. Phys. \textbf{61}, 1--23 (1989).

\bibitem{rubin1970}
V.~C.~Rubin and W.~K.~Ford Jr.,
``Rotation of the Andromeda nebula from a spectroscopic survey
of emission regions,''
Astrophys. J. \textbf{159}, 379--403 (1970).

\bibitem{sparc}
F.~Lelli, S.~S.~McGaugh, and J.~M.~Schombert,
``SPARC: Mass models for 175 disk galaxies with Spitzer
photometry and accurate rotation curves,''
Astron. J. \textbf{152}, 157 (2016).

\bibitem{nfw1997}
J.~F.~Navarro, C.~S.~Frenk, and S.~D.~M.~White,
``A universal density profile from hierarchical clustering,''
Astrophys. J. \textbf{490}, 493--508 (1997).

\bibitem{milgrom1983}
M.~Milgrom,
``A modification of the Newtonian dynamics as a possible
alternative to the hidden mass hypothesis,''
Astrophys. J. \textbf{270}, 365--370 (1983).

\bibitem{alcubierre1994}
M.~Alcubierre,
``The warp drive: hyper-fast travel within general relativity,''
Class. Quantum Grav. \textbf{11}, L73--L77 (1994).

\bibitem{osiris_zenodo}
D.~P.~Davis,
``OSIRIS-CLI: Sovereign quantum discovery framework,''
Zenodo (2025), DOI pending.

\end{thebibliography}

\end{document}
